\section{Step 10 Keep the scalar bases and copy offset}\label{chacha-s10}

The supplied scalar frame protects its incoming return address and three saved registers. Put the immutable input base in \texttt{s0}, working output base in \texttt{s1}, and byte offset zero in \texttt{t0}. This separates algorithm storage from loop progress.

\subsection{Write and inspect}
Replace the body of region \texttt{scalar\_setup}. Retain its \texttt{\# BEGIN} and \texttt{\# END} comments and the supplied labels outside them.

\Needspace{8\baselineskip}
\MARSCodeLabel{MIPS assembly}
\begin{minted}[samepage=true]{mars-idea}
    move $s0, $a0
    move $s1, $a1
    move $t0, $zero
\end{minted}

Use recipe \texttt{s10-1} with the cumulative program through this edit. Inspect its named checkpoint before continuing to the temporary or normal return.

\subsection{Expected observation}
The two bases are retained and the copy offset is zero; the output is still untouched.

\Needspace{9\baselineskip}
\subsection{Check your understanding}
Why keep a base in a saved register?

\begin{itemize}
  \item \textbf{A} Saved registers automatically allocate memory.
  \item \textbf{B} Nested quarter-round calls may change temporary and argument registers.
  \item \textbf{C} A base must fit in thirty-two bits.
\end{itemize}
